\documentclass[11pt]{article}
\usepackage[margin=1in]{geometry}
\usepackage{amsmath,amssymb,amsthm,mathtools}
\usepackage{enumitem}
\usepackage{hyperref}
\usepackage{xcolor}
\usepackage{microtype}
\usepackage{booktabs}
\usepackage{array}

\hypersetup{colorlinks=true, linkcolor=blue!60!black, citecolor=blue!60!black, urlcolor=blue!60!black}

\newtheorem{theorem}{Theorem}[section]
\newtheorem{conjecture}[theorem]{Conjecture}
\newtheorem{lemma}[theorem]{Lemma}
\newtheorem{proposition}[theorem]{Proposition}
\newtheorem{corollary}[theorem]{Corollary}
\newtheorem{definition}[theorem]{Definition}
\newtheorem{remark}[theorem]{Remark}
\newtheorem{problem}[theorem]{Problem}

\newcommand{\one}{\mathbf 1}
\newcommand{\ip}[2]{\left\langle #1,#2\right\rangle}
\newcommand{\Tr}{\operatorname{Tr}}
\newcommand{\norm}[1]{\left\|#1\right\|}
\newcommand{\eps}{\varepsilon}
\newcommand{\pos}[1]{\left(#1\right)_{+}}
\newcommand{\cQ}{\mathcal Q}
\newcommand{\cL}{\mathcal L}

\title{Region II of the Goodman--Style Odd-Cycle Bound:\\
a corrected consolidated proof for $\tfrac12<p<\tfrac23$}
\author{Consolidated working note for T.K., H.Y., and collaborators\\[1mm]
\normalsize Operator and Huber reduction: GPT~Pro, July 11, 2026;\\
\normalsize scalar zone argument and certificates: Claude, July 11, 2026;\\
\normalsize audit, corrections, and consolidation: July 12, 2026}
\date{}

\begin{document}
\maketitle

\begin{abstract}
Let $W$ be a graphon of edge density $p$, and let $m\ge3$ be odd.  This
working note proves the remaining density range
\[
        \frac12<p<\frac23
\]
of the proposed Goodman--style inequality
\[
        t(C_m,W)\ge p^m-p(1-p)^{m-1}.
\]
The proof has two logically separate parts.  First, an operator-theoretic
argument reduces every possible obstruction to one frontier eigenvalue
$\alpha$ of the complement compression and then, through an exact Huber-type
shape elimination, to the explicit scalar inequality
\[
        R_m\le C_m\psi(\xi,\rho).
\]
Second, that scalar inequality is proved on the full admissible parameter
region.  The pinch regime near $p=1/2$, the small-frontier regime, and the
Tur\'an-corner sliver are handled analytically.  Two compact intermediate
regions are covered by finite, exhaustive certificates in exact rational
arithmetic.  Thus the final proof is rigorous but computer-assisted; the
computer is used only for two finite semialgebraic/elementary-function
covering checks, not for floating-point sampling.

This version repairs the supplied July~11 manuscript in four places: it
corrects a reversed comparison in the $m$-maximisation lemma; states and
proves the frontier-gap bound used by the Zone-C certifier; replaces several
unsafe decimal roundings by directed rational bounds; and makes the logical
status of the exact certificates explicit.  The earlier proposed five-atom
stop-loss argument is \emph{not} used here: its scalar variables do not match
the Huber interface, and no valid implication from that lemma to
$R_m\le C_m\psi(\xi,\rho)$ is currently known.  The present document therefore
contains only the proof whose bridge to the graphon problem has been fully
checked.
\end{abstract}

{\small\tableofcontents}

\section{Introduction and main result}\label{sec:intro}

Let $W:[0,1]^2\to[0,1]$ be a graphon with edge density $p=t(K_2,W)$, and let
$m\ge3$ be odd.  The conjecture under study is
\begin{equation}\label{eq:conj}
        t(C_m,W)\ \ge\ g_m(p):=p^m-p(1-p)^{m-1}.
\end{equation}
The consolidated record \texttt{earlier\_consolidated\_note.tex} (July 5, 2026) proves
\eqref{eq:conj}: for $p\le 1/2$ (trivial, $g_m(p)\le0$); for every odd $m$ when
$p\ge2/3$ (the high-density theorem); for all $p$ and all odd $m\le13$
(per-cycle theorems $C_3$--$C_{13}$); and for all $p$, all odd $m$, on several
structured families ($p$-regular; positive-spectrum-safe; rank-one, rank-two and
bipodal complements; equal cliques; two-value multipartite).  The single
remaining case is \emph{Region II}:
\[
        \tfrac12<p<\tfrac23,
        \qquad\text{i.e.}\qquad
        q:=1-p\in\left(\tfrac13,\tfrac12\right),
        \qquad m\ge15\ \text{odd}.
\]

A July 11 working derivation by GPT~Pro reduces Region II rigorously to a single \emph{scalar} inequality --- the
\emph{Huber payment inequality} --- in the three variables $(q,\alpha,m)$, where
$\alpha$ is the unique eigenvalue of the complement compression above $q$
(the ``frontier eigenvalue''); all other spectral, eigenfunction-shape and
coupling data are eliminated exactly.  Section~\ref{sec:reduction} restates that
reduction, self-contained and with complete proofs, since everything that
follows rests on it.

The present note proves the scalar inequality.

\begin{theorem}[Scalar Huber inequality]\label{thm:scalar}
Let $q\in(\tfrac13,\tfrac12)$, let $\alpha\in\bigl(q,\,r(q)\bigr]$ with
$r(q)=\frac{\sqrt{q^2+4q}-q}2$, and let $m\ge15$ be odd.  With the notation of
Definition~\ref{def:huber} (all quantities explicit functions of
$(q,\alpha,m)$),
\[
        R_m\ \le\ C_m\,\psi(\xi,\rho).
\]
\end{theorem}

\begin{corollary}[Region II]\label{cor:region2}
For every graphon $W$ with $\tfrac12<p<\tfrac23$ and every odd $m\ge3$,
\[
        t(C_m,W)\ \ge\ p^m-p(1-p)^{m-1}.
\]
\end{corollary}

\begin{corollary}[The full conjecture]\label{cor:full}
For every graphon $W$ and every odd $m\ge3$,
$\ t(C_m,W)\ \ge\ p^m-p(1-p)^{m-1}$, $p=t(K_2,W)$.
\end{corollary}

\paragraph{Proof architecture.}
The scalar inequality is proved through the dual formula for the Huber
envelope $\psi$.  Zones A and B use $\lambda=1$; Zone C uses either
$\lambda=1$ or $\lambda=2\rho_0\xi$, where $0<\rho_0\le\rho$ is an
explicit lower bound and the latter choice is made only when it lies in
$[0,1]$.
The admissible domain is covered by three zones: the pinch zone A
($e=1-2\alpha\le1/60$, $\xi\ge1$), zone B
($e\ge1/60$, $\xi\ge1$), and zone C ($\xi\le1$).  Zone A and the
small-$e$ part of Zone C are elementary analytic estimates.  Zone B is an
analytic reduction to a two-variable inequality followed by a finite exact
rational box certificate.  The moderate-$e$ part of Zone C is an analytic
three-geometric reduction followed by a second finite exact rational
certificate; its very thin Tur\'an-corner remainder is again analytic.

\subsection*{Provenance and verification status}

\begin{center}
\small
\begin{tabular}{>{\raggedright\arraybackslash}p{0.27\linewidth}
                  >{\raggedright\arraybackslash}p{0.20\linewidth}
                  >{\raggedright\arraybackslash}p{0.45\linewidth}}
\toprule
Result & Source & Status in this version \\
\midrule
One-frontier and Huber reduction (\S\ref{sec:reduction})
& GPT~Pro, July 11
& Pencil-and-paper proof; the endpoint issue in the forced-variance argument
and the $K=0$ shape case are stated explicitly. \\
Chart and scalar certificates (\S\ref{sec:coords})
& Claude, July 11
& Pencil-and-paper identities; exact symbolic regression checks supplied. \\
Zone A (\S\ref{sec:zoneA})
& Claude, corrected here
& Pencil-and-paper proof with all decimal constants interpreted as exact
rationals and rounded in the safe direction. \\
Zone B (\S\ref{sec:zoneB})
& Claude, corrected here
& Analytic $m$-maximisation plus the exact-rational verifier
\texttt{zoneB\_certifier.py}; the reversed comparison in the manuscript is fixed. \\
Zone C (\S\ref{sec:zoneC})
& Claude, corrected here
& Analytic for $e\le1/60$ and in the Tur\'an corner; analytic reduction plus
\texttt{zoneC\_certifier.py} on the compact middle range.  The missing bound
$f\ge d+\delta$ is now proved and used. \\
Numerical sweep (\S\ref{sec:verification})
& Companion checker
& Regression test only; no theorem depends on sampled floating-point values. \\
\bottomrule
\end{tabular}
\end{center}

\begin{remark}[Why the stop-loss note is not spliced into this proof]
A separate five-atom stop-loss inequality was derived during the audit.  Its
parameters use a different quantity
$(q-\alpha^2)/(1-\alpha)$, whereas the Huber reduction uses
$L=\sqrt{pq-\alpha^2}$.  The tempting domination that would bridge the two
statements is false in general (already for admissible parameters at large
$m$).  Accordingly, that lemma is not invoked anywhere below.  This remark is
included to make the dependency graph unambiguous for later paper writing.
\end{remark}

\section{The one-frontier reduction, recapped with proofs}\label{sec:reduction}

Everything in this section is due to the July 11 baton note; we restate it so
that the present document is self-contained, and we flag two small points of
care (Remarks~\ref{rmk:budget-sign} and \ref{rmk:knegative}) that the original
exposition leaves implicit.

\subsection{Operator set-up and the Hilbert--Schmidt budget}

Let $U=1-W$, let $T_U$ be the integral operator with kernel $U$, and decompose
$L^2([0,1])=\mathbb R\one\oplus\one^\perp$.  In block form,
\begin{equation}\label{eq:blockTU}
        T_U=
        \begin{pmatrix}
        q&g^*\\
        g&A
        \end{pmatrix},
        \qquad
        g=T_U\one-q\one\in\one^\perp,
        \qquad
        A=P_{\one^\perp}T_UP_{\one^\perp}.
\end{equation}
Since $0\le U\le1$ pointwise, $\int U^2\le\int U=q$, so
$\|T_U\|_{\mathrm{HS}}^2\le q$; expanding in blocks,
\begin{equation}\label{eq:HS-budget}
        q^2+2\|g\|_2^2+\Tr(A^2)\ \le\ q,
        \qquad\text{i.e.}\qquad
        \Tr(A^2)+2\|g\|_2^2\ \le\ pq.
\end{equation}

\subsection{The one-sided spectral shift}

\begin{proposition}[One-sided spectral shift]\label{prop:shift}
For every odd $m\ge3$,
\begin{equation}\label{eq:one-sided-shift}
        t(C_m,W)=p^m-\Tr(A^m)+m\,[z^m]\,\cL_W(z),
        \qquad
        \cL_W(z)=-\log\Bigl(1-\frac{z^2\ip{g}{(I+zA)^{-1}g}}{1-pz}\Bigr),
\end{equation}
as an identity of formal power series coefficients.
\end{proposition}

\begin{proof}
For finite-rank step graphons, $t(C_m,W)=\Tr(T_W^m)$ with $T_W=J-T_U$,
$Jf=\ip f\one\one$.  Conjugating by the sign flip on $\one^\perp$ shows
$\Tr(T_W^m)=\Tr(M_W^m)$ with
$M_W=\bigl(\begin{smallmatrix}p&g^*\\ g&-A\end{smallmatrix}\bigr)$.
Schur's determinant formula gives
\[
        \det(I-zM_W)=\det(I+zA)\Bigl(1-pz-z^2\ip{g}{(I+zA)^{-1}g}\Bigr).
\]
Taking $-\log$ and using $-\log\det(I-zT)=\sum_{r\ge1}z^r\Tr(T^r)/r$,
\[
        \sum_{r\ge1}\frac{z^r}{r}\Tr(M_W^r)
        =-\log\det(I+zA)-\log(1-pz)+\cL_W(z),
\]
because $1-pz-z^2\ip{g}{(I+zA)^{-1}g}
=(1-pz)\bigl(1-\frac{z^2\ip{g}{(I+zA)^{-1}g}}{1-pz}\bigr)$.
Comparing the coefficient of $z^m$ and multiplying by $m$ yields
\eqref{eq:one-sided-shift}, since $\Tr((-A)^m)=-\Tr(A^m)$ for odd $m$.
For general graphons, let $U_n$ be the conditional expectation of $U$ on the
dyadic grid of mesh $2^{-n}$; then $0\le U_n\le1$, $\int U_n=q$, and $U_n\to U$
in $L^2$.  Each side of \eqref{eq:one-sided-shift} is continuous along this
approximation: $t(C_m,\cdot)$ and $\Tr(A^m)$ by telescoping and
Schatten--H\"older ($m\ge3\ge2$), and $m[z^m]\cL_W$ because it is a polynomial
in $p$ and the finitely many moments $\ip g{A^jg}$, $j\le m-2$, each continuous
under $(g_n,A_n)\to(g,A)$ in $L^2\times\mathrm{HS}$.
\end{proof}

\begin{lemma}[Compression bound and the no-frontier case]\label{lem:no-frontier}
The compression $A$ satisfies $\|A\|\le\tfrac12$.  Moreover, if
$\lambda_{\max}(A)\le q$, then for every odd $m\ge3$,
\[
        t(C_m,W)\ge p^m-pq^{m-1}.
\]
\end{lemma}

\begin{proof}
For a real $h\perp\one$, write $h=h^+-h^-$ and
$r=\int h^+=\int h^-=\tfrac12\|h\|_1$.  Since $0\le U\le1$,
\[
 -2r^2\le\langle h,T_Uh\rangle\le2r^2
 \le\frac12\|h\|_2^2.
\]
Taking the supremum over unit $h\in\one^\perp$ gives $\|A\|\le1/2<p$.

As in the proof of Proposition~\ref{prop:master-defect}, the formal series
\[
 u(z)=\frac{z^2\langle g,(I+zA)^{-1}g\rangle}{1-pz}
\]
has nonnegative coefficients: for an eigenvalue $\nu\in[-1/2,1/2]$ of $A$,
the coefficient of $z^n$, $n\ge2$, contributed by its coupling is
\[
 \frac{p^{n-1}-(-\nu)^{n-1}}{p+\nu}\ge0.
\]
Consequently $[z^m]\mathcal L_W(z)\ge0$.  If
$\lambda_{\max}(A)\le q$, then for every eigenvalue $\nu$ of $A$ and odd
$m\ge3$,
\[
        \nu^m\le q^{m-2}\nu^2.
\]
Together with $\operatorname{Tr}(A^2)\le pq$, from
\eqref{eq:HS-budget}, this yields
$\operatorname{Tr}(A^m)\le pq^{m-1}$.  Proposition~\ref{prop:shift} now gives
the claim.
\end{proof}

\subsection{One frontier, forced variance, and the frontier ceiling}

Fix $q\in(1/3,1/2)$.  By Lemma~\ref{lem:no-frontier}, only the case in which
$A$ has an eigenvalue $\alpha>q$ remains.  Fix a corresponding unit
eigenfunction $\phi\perp\one$.

\begin{lemma}[One frontier]\label{lem:one-frontier}
$A$ has at most one eigenvalue $>q$, and every other eigenvalue lies in
$[-L,L]$, where
\[
        L:=\sqrt{pq-\alpha^2}\,,\qquad 0\le L<q<\alpha<\tfrac12 .
\]
\end{lemma}

\begin{proof}
Two eigenvalues $>q$ would give $\Tr(A^2)>2q^2>q(1-q)\ge\Tr(A^2)$ for
$q>1/3$, by \eqref{eq:HS-budget}: contradiction.  By \eqref{eq:HS-budget}
again, the total square mass of the non-frontier spectrum is at most
$pq-\alpha^2=L^2$, so each non-frontier eigenvalue has modulus at most $L$.
Finally $L^2=pq-\alpha^2<pq-q^2=q(1-2q)<q^2$ for $q>\tfrac13$, so $L<q$; and
$\alpha\le\lambda_{\max}(A)\le\|A\|\le\|T_U\|_{\mathrm{HS}}\le\sqrt q<1$,
while $\alpha<\tfrac12$ follows from Lemma~\ref{lem:forced-variance} below.
\end{proof}

\begin{lemma}[Forced variance and the frontier ceiling]\label{lem:forced-variance}
Write $a=\int|\phi|$.  Then $a^2\ge2\alpha$, and
\begin{equation}\label{eq:forced-variance}
        \|g\|_2^2\ \ge\ \frac{\alpha(\alpha-q)^2}{2(1-2\alpha)} .
\end{equation}
Consequently
\begin{equation}\label{eq:frontier-ceiling}
        \alpha\ \le\ r(q):=\frac{\sqrt{q^2+4q}-q}{2},
\end{equation}
i.e.\ $\alpha^2+q\alpha-q\le0$.  (In particular $\alpha\le r(q)<r(1/2)=1/2$.)
\end{lemma}

\begin{proof}
Split $\phi=\phi^+-\phi^-$.  Since $U\ge0$,
$\alpha=\ip{\phi}{T_U\phi}\le\ip{\phi^+}{T_U\phi^+}+\ip{\phi^-}{T_U\phi^-}$,
and since $U\le1$, $\ip{\phi^\pm}{T_U\phi^\pm}\le(\int\phi^\pm)^2=(a/2)^2$
(both signs have integral $a/2$ because $\int\phi=0$).  Hence $\alpha\le a^2/2$.

Let $f=|\phi|=a\one+h$, $h\perp\one$, $\|h\|_2^2=1-a^2$.  Since $U\ge0$,
$4\ip{\phi^+}{T_U\phi^-}=\ip f{T_Uf}-\ip\phi{T_U\phi}\ge0$, so
$\ip f{T_Uf}\ge\alpha$.  Expanding via \eqref{eq:blockTU} and
$\ip h{Ah}\le\alpha\|h\|_2^2$,
\[
        \alpha\le qa^2+2a\ip gh+\alpha(1-a^2)
        \le qa^2+2a\|g\|_2\sqrt{1-a^2}+\alpha(1-a^2),
\]
whence $(\alpha-q)a^2\le2a\|g\|_2\sqrt{1-a^2}$.  Here $a^2<1$:
if $a^2=1$, then the right-hand side vanishes while $\alpha>q$, a
contradiction.  Since $a^2\ge2\alpha$, this also gives $\alpha<1/2$.  We may
therefore divide by $1-a^2$ and obtain
\[
        \|g\|_2^2\ge\frac{(\alpha-q)^2a^2}{4(1-a^2)}
        \ge\frac{\alpha(\alpha-q)^2}{2(1-2\alpha)},
\]
using the monotonicity of $t\mapsto t/(1-t)$.
Combining \eqref{eq:forced-variance} with $\alpha^2+2\|g\|_2^2\le pq$ (from
\eqref{eq:HS-budget}, $\Tr(A^2)\ge\alpha^2$) and clearing the positive
denominator $1-2\alpha$ gives
$(1-\alpha-q)(\alpha^2+q\alpha-q)\le0$; since $\alpha+q<1$, the ceiling
\eqref{eq:frontier-ceiling} follows.
\end{proof}

\begin{remark}\label{rmk:budget-sign}
Two consequences of \eqref{eq:HS-budget} are used repeatedly below:
$\alpha^2+2\|g\|_2^2\le pq$, i.e.\ $\|g\|_2^2\le L^2/2$; and the non-frontier
square mass is at most $L^2-2\|g\|_2^2\ge0$.  The nonnegativity of the latter
(automatic, by the former) is what makes the square-mass bookkeeping in
Proposition~\ref{prop:master-defect} legitimate in all configurations.
\end{remark}

\subsection{The master defect estimate}

Decompose the degree fluctuation along the frontier:
\begin{equation}\label{eq:g-decomp}
        g=c\phi+g_s,\qquad c=\ip g\phi,\qquad g_s\perp\phi .
\end{equation}
For odd $m$ define on $(-p,p)$
\begin{equation}\label{eq:kdef}
        k_m(\lambda)=\frac{p^{m-1}-\lambda^{m-1}}{p+\lambda},
\end{equation}
and set
\begin{equation}\label{eq:AB-def}
        A_m=2L^{m-2}+m\,k_m(\alpha),
        \qquad
        B_m=2L^{m-2}+m\,k_m(L),
        \qquad
        R_m=\alpha^m+L^m-pq^{m-1}.
\end{equation}

\begin{remark}\label{rmk:knegative}
$k_m$ is positive on $(-p,p)$, strictly decreasing on $[0,p)$ (from
$k_m'(\lambda)(p+\lambda)=-(m-1)\lambda^{m-2}-k_m(\lambda)$), and for
$\lambda\in[-L,0]$ one has
$k_m(\lambda)=\sum_{j=0}^{m-2}(-\lambda)^jp^{m-2-j}\ge p^{m-2}=k_m(0)\ge k_m(L)$.
Hence $k_m(\lambda)\ge k_m(L)$ for \emph{all} $\lambda\in[-L,L]$, and
$0<A_m<B_m$.
\end{remark}

\begin{proposition}[Master defect]\label{prop:master-defect}
Suppose $\|g_s\|_2^2\ge\tau$.  Then
\begin{equation}\label{eq:master-defect}
        t(C_m,W)-\bigl(p^m-pq^{m-1}\bigr)\ \ge\ -R_m+A_mc^2+B_m\tau .
\end{equation}
\end{proposition}

\begin{proof}
We bound the two non-baseline terms of \eqref{eq:one-sided-shift}.

\emph{(i) The trace term.}  By Lemma~\ref{lem:one-frontier} and
Remark~\ref{rmk:budget-sign}, the non-frontier eigenvalues $\lambda_i$ lie in
$[-L,L]$ and satisfy $\sum_i\lambda_i^2\le L^2-2\|g\|_2^2$.  For odd $m$,
\[
        \Tr(A^m)=\alpha^m+\sum_i\lambda_i^m
        \le\alpha^m+\sum_{\lambda_i>0}\lambda_i^m
        \le\alpha^m+L^{m-2}\!\!\sum_{\lambda_i>0}\!\lambda_i^2
        \le\alpha^m+L^m-2L^{m-2}\bigl(c^2+\|g_s\|_2^2\bigr).
\]

\emph{(ii) The shift term.}  Write
$u(z)=\frac{z^2\ip{g}{(I+zA)^{-1}g}}{1-pz}$.  Expanding $g$ in the spectral
resolution of $A$ (eigenvalues $\nu_i\in[-L,L]\cup\{\alpha\}$, couplings
$c_i^2$),
\[
        [z^n]\,u(z)=\sum_i c_i^2\,
        \frac{p^{n-1}-(-\nu_i)^{n-1}}{p+\nu_i}\ \ge\ 0
        \qquad(n\ge2),
\]
since $|\nu_i|<p$.  Thus $u$ has nonnegative coefficients, hence so does
$\cL_W=\sum_{r\ge1}u^r/r$, and
\[
        m[z^m]\cL_W\ \ge\ m[z^m]u
        =m\sum_ic_i^2k_m(\nu_i)
        \ \ge\ m\,k_m(\alpha)\,c^2+m\,k_m(L)\,\|g_s\|_2^2,
\]
using $(-\nu)^{m-1}=\nu^{m-1}$ ($m-1$ even), Remark~\ref{rmk:knegative}, and
$\sum_{i\,\mathrm{safe}}c_i^2=\|g_s\|_2^2$.  Combining (i), (ii) with
\eqref{eq:one-sided-shift} and $\|g_s\|_2^2\ge\tau$ gives
\eqref{eq:master-defect}.
\end{proof}

\subsection{Eigenfunction geometry: the two coupling channels}

Set
\begin{equation}\label{eq:shape-defs}
        z=\Bigl(\int|\phi|\Bigr)^2,\quad
        b=\ip{|\phi|}{\phi}\ \ (\ge0\ \text{after}\ \phi\mapsto\pm\phi),\quad
        k=|\phi|-\sqrt z\,\one-b\phi,\quad K=\|k\|_2^2 ,
\end{equation}
so that $z+b^2+K=1$ and, by Lemma~\ref{lem:forced-variance}, $z\ge2\alpha$.
Write $h=z-2\alpha\ge0$ and $e=1-2\alpha>0$, so $h+b^2+K=e$.

\begin{lemma}[Direct channel]\label{lem:direct}
$\displaystyle c\ \le\ u_c:=\frac{h-2\alpha b}{2\sqrt z}$.
\end{lemma}

\begin{proof}
$T_U\phi=c\one+\alpha\phi$ by \eqref{eq:blockTU}; also
$T_U\phi\le T_U\phi^+\le\int\phi^+=\sqrt z/2$ pointwise ($U\le1$,
$\phi\le\phi^+$).  Pairing with $\phi^+\ge0$, whose squared norm is
$\ip{\phi^+}{\phi^+}=\tfrac{1+b}2$ (from
$\|\phi^+\|^2-\|\phi^-\|^2=b$, $\|\phi^+\|^2+\|\phi^-\|^2=1$):
\[
        \frac{c\sqrt z}{2}+\alpha\,\frac{1+b}{2}
        =\ip{\phi^+}{T_U\phi}
        \le\frac{\sqrt z}{2}\int\phi^+=\frac z4 ,
\]
which rearranges to the claim using $z-2\alpha=h$.
\end{proof}

\begin{lemma}[Safe channel]\label{lem:safe}
With $H=\dfrac{(\alpha-q)z+(\alpha-L)K}{2\sqrt z}$,
\[
        bc+\ip{g_s}{k}\ \ge\ H,
        \qquad\text{hence}\qquad
        \|g_s\|_2^2\ \ge\ \frac{\pos{H-bc}^2}{K}\quad(K>0).
\]
\end{lemma}

\begin{proof}
$U\ge0$ gives $\ip{|\phi|}{T_U|\phi|}\ge\ip{\phi}{T_U\phi}=\alpha$.  Expanding
the left side via $|\phi|=\sqrt z\one+b\phi+k$ and \eqref{eq:blockTU}
($\ip gk=\ip{g_s}k$ because $k\perp\phi$; $\ip{\phi}{Ak}=0$):
\[
        qz+2\sqrt z\,(bc+\ip{g_s}k)+\alpha b^2+\ip k{Ak}\ \ge\ \alpha
        =\alpha(z+b^2+K).
\]
Since $k\perp\one,\phi$ lies in the safe subspace, $\ip k{Ak}\le LK$
(Lemma~\ref{lem:one-frontier}); rearranging gives the first claim, and
Cauchy--Schwarz $\ip{g_s}k\le\|g_s\|_2\sqrt K$ the second.
\end{proof}

\subsection{Exact Huber elimination}

\begin{definition}\label{def:huber}
For admissible $(q,\alpha,m)$ put $d=\alpha-q$, $f=\alpha-L$, $e=1-2\alpha$, and
\begin{equation}\label{eq:Cxirho}
        C_m=\frac{B_mf\sqrt{2\alpha}\,e^2}{4\alpha^2},
        \qquad
        \xi=\frac{4\alpha^2d}{e^2},
        \qquad
        \rho=\frac{A_m}{B_m}\,\frac{\sqrt\alpha}{2\sqrt2\,f},
\end{equation}
\begin{equation}\label{eq:psi}
        \psi(\xi,\rho)=\min_{0\le v\le1}\bigl\{\rho v^2+\pos{\xi-v+v^2}\bigr\}.
\end{equation}
\end{definition}

\begin{theorem}[Exact Huber shape elimination; GPT~Pro]\label{thm:huber}
For every Region-II graphon configuration with frontier data
$(q,\alpha,c,g_s,\phi)$,
\[
        A_mc^2+B_m\|g_s\|_2^2\ \ge\ C_m\,\psi(\xi,\rho).
\]
Consequently, by Proposition~\ref{prop:master-defect}, Conjecture \eqref{eq:conj}
holds in Region II as soon as the scalar inequality
\begin{equation}\label{eq:scalar-target}
        R_m\ \le\ C_m\,\psi(\xi,\rho)
\end{equation}
holds for all $q\in(\tfrac13,\tfrac12)$, $\alpha\in(q,r(q)]$ and odd $m\ge15$.
\end{theorem}

\begin{proof}
Write $\cQ=\min\{A_mc^2+B_mK^{-1}\pos{H-bc}^2:\ c\le u_c\}$ over the admissible
shape data; by Lemmas~\ref{lem:direct} and~\ref{lem:safe} the left side
dominates $\cQ$ pointwise, so it suffices to prove
$\cQ\ge C_m\psi(\xi,\rho)$.  (Degenerate case $K=0$: then $k=0$ and
Lemma~\ref{lem:safe} gives
$bc\ge H=d\sqrt z/2\ge d\sqrt{2\alpha}/2>0$.  Hence $c>0$, so
$u_c>0$ and the budget \eqref{eq:bc-budget} below applies.  With
$v=2c\sqrt{2\alpha}/e$ it gives
$b\le\beta(c)\le e(1-v)/(2\alpha)$, and therefore
\[
 \frac{d\sqrt{2\alpha}}2
 \le bc\le c\beta(c)
 \le\frac{e^2v(1-v)}{4\alpha\sqrt{2\alpha}}.
\]
Equivalently $\xi\le v-v^2$.  Thus the positive-part term in
\eqref{eq:psi} vanishes at this $v$, and
$A_mc^2=C_m\rho v^2\ge C_m\psi$.  We may henceforth assume $K>0$.)

\emph{Case $u_c\le0$.}  Every admissible $c$ is $\le0$; since $b\ge0$,
$H-bc\ge H$, so $\cQ\ge B_mH^2/K$.  By AM--GM,
\[
        \frac{H^2}{K}=\frac{\bigl((\alpha-q)z+(\alpha-L)K\bigr)^2}{4zK}
        \ \ge\ (\alpha-q)(\alpha-L)=df ,
\]
and $C_m\psi\le C_m\xi=B_mf\sqrt{2\alpha}\,d\le B_mfd$ (take $v=0$ in
\eqref{eq:psi}; $\sqrt{2\alpha}<1$), so $\cQ\ge C_m\psi$.

\emph{Case $u_c>0$.}  A minimizer may be taken with $c\ge0$: replacing $c<0$ by
$0$ decreases both terms (for the second, $H-bc\ge H\ge H-b\cdot0$ when $c\le0$,
$b\ge0$).  The constraint $c\le u_c$, i.e.\
$2c\sqrt z+2\alpha b\le h=e-b^2-K\le e-b^2$, together with $z\ge2\alpha$, gives
\begin{equation}\label{eq:bc-budget}
        b^2+2\alpha b+2c\sqrt{2\alpha}\ \le\ e ,
\end{equation}
so $0\le c\le e/(2\sqrt{2\alpha})$ and
$b\le\beta(c):=\sqrt{\alpha^2+e-2c\sqrt{2\alpha}}-\alpha
\le\frac{e-2c\sqrt{2\alpha}}{2\alpha}$.
Also $z\le1$ and $z\ge2\alpha$ give
$H\ge\frac{d\sqrt{2\alpha}}2+\frac f2K$.
For fixed $c$ and any $K>0$, AM--GM yields
$\frac{(w+fK/2)^2}{K}\ge2fw$ for $w\ge0$, hence
\[
        \cQ\ \ge\ \min_{0\le c\le e/(2\sqrt{2\alpha})}
        \Bigl\{A_mc^2+2B_mf\,\pos{\tfrac{d\sqrt{2\alpha}}2-c\beta(c)}\Bigr\}.
\]
Substituting $v=2c\sqrt{2\alpha}/e\in[0,1]$ and using
$c\beta(c)\le\frac{e^2v(1-v)}{4\alpha\sqrt{2\alpha}}$:
\[
        \frac{d\sqrt{2\alpha}}2-c\beta(c)
        \ \ge\ \frac{\sqrt{2\alpha}\,e^2}{8\alpha^2}\,(\xi-v+v^2),
        \qquad
        A_mc^2=\frac{A_me^2}{8\alpha}v^2=C_m\rho v^2 ,
\]
and $2B_mf\cdot\frac{\sqrt{2\alpha}e^2}{8\alpha^2}=C_m$, so
$\cQ\ge C_m\min_{0\le v\le1}\{\rho v^2+\pos{\xi-v+v^2}\}=C_m\psi(\xi,\rho)$.
\end{proof}

\begin{proposition}[Forms of $\psi$; GPT~Pro]\label{prop:psi-forms}
Let $\xi_c(\rho)=\frac{2\rho+1}{4(\rho+1)^2}$ and, for $0\le\xi<\frac14$,
$v_-(\xi)=\frac{1-\sqrt{1-4\xi}}2$.  Then
\[
        \psi(\xi,\rho)=
        \begin{cases}
        \rho\,v_-(\xi)^2,&0\le\xi<\xi_c(\rho),\\[1mm]
        \xi-\dfrac1{4(1+\rho)},&\xi\ge\xi_c(\rho),
        \end{cases}
        \qquad\text{and}\qquad
        \psi(\xi,\rho)=\max_{0\le\lambda\le1}
        \Bigl(\lambda\xi-\frac{\lambda^2}{4(\rho+\lambda)}\Bigr).
\]
\end{proposition}

\begin{proof}
On $\{v:\xi-v+v^2\le0\}$ (nonempty iff $\xi\le1/4$) the objective is $\rho v^2$,
minimized at $v_-(\xi)$; elsewhere it is $(\rho+1)v^2-v+\xi$ with unconstrained
minimizer $v_*=\frac1{2(1+\rho)}\in(0,1)$ and value $\xi-\frac1{4(1+\rho)}$;
the branch transition $v_*=v_-(\xi)$ occurs exactly at $\xi=\xi_c(\rho)$.
For the dual form, write $\pos t=\max_{0\le\lambda\le1}\lambda t$, exchange
$\min_v\max_\lambda$ (convex--concave on compact intervals), and minimize
$(\rho+\lambda)v^2-\lambda v+\lambda\xi$ over $v\in[0,1]$: the minimizer
$\lambda/(2(\rho+\lambda))\in[0,1]$ gives the stated value.
\end{proof}

The dual form furnishes the two \emph{certificates} used throughout:
\begin{align}
\lambda=1:&\qquad
        C_m\psi\ \ge\ C_m\Bigl(\xi-\frac1{4(1+\rho)}\Bigr)
        =\sqrt{2\alpha}\,B_mf\Bigl(d-\frac{e^2}{16\alpha^2(1+\rho)}\Bigr),
        \label{eq:cert-II}\\
\lambda=2\rho\xi\ (\le1):&\qquad
        C_m\psi\ \ge\ C_m\,\rho\,\xi^2\,\frac{1+4\xi}{1+2\xi}
        =\frac{2\alpha^3A_md^2}{e^2}\cdot\frac{1+4\xi}{1+2\xi},
        \label{eq:cert-I}
\end{align}
where we used $C_m\xi=\sqrt{2\alpha}B_mfd$ and
$C_m\rho\xi^2=2\alpha^3A_md^2/e^2$ (both exact, from
\eqref{eq:Cxirho}).

\section{The scalar problem in frontier--gap coordinates}\label{sec:coords}

\subsection{The admissible domain and the \texorpdfstring{$(e,\kappa)$}{(e, kappa)} chart}

\begin{definition}[Admissible domain]\label{def:domain}
$\mathcal D$ is the set of triples $(q,\alpha,m)$ with $q\in(\tfrac13,\tfrac12)$,
$\alpha\in(q,r(q)]$, $m\ge15$ odd.  Throughout, $p=1-q$, $d=\alpha-q$,
$e=1-2\alpha$, $L=\sqrt{pq-\alpha^2}$, $f=\alpha-L$, and we use the ratios
\[
        x=\frac\alpha p,\quad
        \tau=\frac q\alpha,\quad
        y=\frac Lp,\quad
        \ell=\frac L\alpha,\quad
        u=1-x=\frac{d+e}{p},\quad
        T=me,\quad
        \kappa=\frac de .
\]
\end{definition}

\begin{lemma}[The $(e,\kappa)$ chart]\label{lem:chart}
The map $(q,\alpha)\mapsto(e,\kappa)$ identifies
$\{(q,\alpha):q\in(\tfrac13,\tfrac12),\ \alpha\in(q,r(q)]\}$ with
\[
        \Bigl\{(e,\kappa):\ 0<e<\tfrac13,\quad
        0<\kappa\le\kappa_{\max}(e):=\frac{1-e}{1+e},\quad
        \kappa<\kappa_q(e):=\frac{1-3e}{6e}\Bigr\},
\]
via $\alpha=\frac{1-e}2$, $q=\alpha-\kappa e$, $p=\alpha+(1+\kappa)e$.  On this
chart,
\[
        L^2=\alpha e-d(d+e),
        \qquad
        \xi=\frac{(1-e)^2\kappa}{e},
        \qquad
        \frac1x=1+\frac{2(1+\kappa)e}{1-e},
        \qquad
        1-\tau=\frac{2\kappa e}{1-e}.
\]
\end{lemma}

\begin{proof}
$e=1-2\alpha\in(0,\tfrac13)$ iff $\alpha\in(\tfrac13,\tfrac12)$, which holds
since $q>\tfrac13$ and $\alpha\le r(q)<r(\tfrac12)=\tfrac12$.  Next,
$q>\tfrac13$ iff $\kappa e<\alpha-\tfrac13=\frac{1-3e}6$, i.e.\
$\kappa<\kappa_q(e)$.  For the ceiling: a direct computation gives
\[
        \alpha^2+q\alpha-q
        =\alpha^2+q(\alpha-1)
        =\frac{(1-e)^2}4-\Bigl(\frac{1-e}2-\kappa e\Bigr)\frac{1+e}2
        =\frac e2\bigl(\kappa(1+e)-(1-e)\bigr),
\]
so $\alpha\le r(q)$ (i.e.\ $\alpha^2+q\alpha-q\le0$) iff
$\kappa\le\kappa_{\max}(e)$.  Also
$L^2=pq-\alpha^2=(\alpha+d+e)(\alpha-d)-\alpha^2=\alpha e-d(d+e)$,
$\xi=4\alpha^2d/e^2=(1-e)^2\kappa/e$, and the last two identities are
immediate from $p=\alpha+(1+\kappa)e$, $q=\alpha-\kappa e$,
$\alpha=\frac{1-e}2$.
\end{proof}

\begin{lemma}[Elementary facts on $\mathcal D$]\label{lem:elementary}
On $\mathcal D$: $0\le L<q<\alpha<\tfrac12<p$; $y<\tfrac12$;
$\ell\le\sqrt{e/\alpha}=\sqrt{2e/(1-e)}$; and for every integer $n\ge2$ and
$0<u<1$,
\begin{equation}\label{eq:E-a}
        (1-u)^n\ \le\ \frac1{1+nu+\binom n2u^2},
        \qquad
        (1-u)^n\ \ge\ 1-nu .
\end{equation}
Moreover the defect has the exact form
\begin{equation}\label{eq:defect-form}
        R_m=\alpha^{m-1}\bigl(\alpha-p\,\tau^{m-1}\bigr)+L^m .
\end{equation}
\end{lemma}

\begin{proof}
$L<q$: $L^2=pq-\alpha^2<pq-q^2=q(1-2q)<q^2$ for $q>\tfrac13$.
$y<\tfrac12$: with $s=q/p\in(\tfrac12,1)$ and $x>s$ one has
$y^2=pq/p^2-x^2=s-x^2<s-s^2\le\tfrac14$.
$\ell^2=L^2/\alpha^2\le\alpha e/\alpha^2=e/\alpha$.
\eqref{eq:E-a}: the binomial series
$(1-u)^{-n}=\sum_k\binom{n+k-1}ku^k$ has nonnegative terms, so
$(1-u)^{-n}\ge1+nu+\binom{n+1}2u^2\ge1+nu+\binom n2u^2$; Bernoulli is
standard.  \eqref{eq:defect-form}: $q=\alpha\tau$ and
$\alpha^m+L^m-pq^{m-1}
=\alpha^{m-1}\alpha-p\alpha^{m-1}\tau^{m-1}+L^m$.
\end{proof}

\begin{lemma}[Frontier-gap inequalities]\label{lem:frontier-gap}
Put
\[
        \delta:=\alpha-\frac13>0.
\]
Then, throughout the admissible domain,
\begin{equation}\label{eq:frontier-gap}
        d+L\le\frac13,
        \qquad q-L\ge\delta,
        \qquad f=\alpha-L\ge d+\delta.
\end{equation}
\end{lemma}

\begin{proof}
Since $q=\alpha-d=\frac13+\delta-d$ and
$L^2=pq-\alpha^2$, direct expansion gives
\[
 \left(\frac13-d\right)^2-L^2
 =2(d^2-d\delta+\delta^2)+\frac{\delta-d}{3}\ge0.
\]
The domain condition $q>1/3$ is exactly $0<d<\delta$, so
$1/3-d>0$.  Hence $L\le1/3-d$, proving the first inequality.  The other two
follow from
\[
 q-L=\alpha-(d+L)\ge\alpha-\frac13=\delta,
 \qquad
 f=d+(q-L).
\]
\end{proof}

\subsection{The three zones}

Theorem~\ref{thm:scalar} is proved by combining three lemmas whose hypotheses
cover $\mathcal D$ (recall the two certificates \eqref{eq:cert-II},
\eqref{eq:cert-I}):

\begin{itemize}[leftmargin=2em]
\item \textbf{Zone A} (\S\ref{sec:zoneA}): $e\le\tfrac1{60}$ and $\xi\ge1$;
certificate $\lambda=1$.  This zone contains the \emph{pinch}
$q\uparrow\tfrac12$, where the inequality is asymptotically tight (relative
margin $\Theta(1/m)$); its proof is the sharpest.
\item \textbf{Zone B} (\S\ref{sec:zoneB}): $e\ge\tfrac1{60}$ and $\xi\ge1$;
certificate $\lambda=1$.  (Nonempty only for $e<0.2033$: $\xi\ge1$ requires
$\kappa\ge e/(1-e)^2$, while the domain requires $\kappa<\kappa_q(e)$; the two
are compatible iff $6e^2<(1-e)^2(1-3e)$, whose left-minus-right side is
strictly increasing in $e$ and changes sign in $(0.2032,\,2033/10000)$ ---
exact rational evaluation at the endpoints.)
\item \textbf{Zone C} (\S\ref{sec:zoneC}): $\xi\le1$ (any $e$); the
admissible dual choice is $\lambda=2\rho_0\xi$ when this is at most $1$,
and $\lambda=1$ otherwise, for an explicit $\rho_0\le\rho$.  This zone
contains the ultra-thin frontier scale
and the \emph{Tur\'an corner} $e\uparrow\tfrac13$, $\kappa\downarrow0$, where
$R_m\to0$ and the constraint $q>\tfrac13$ is essential.
\end{itemize}

\section{Zone A: \texorpdfstring{$e\le1/60$, $\xi\ge1$}{e <= 1/60, xi >= 1}}\label{sec:zoneA}

\begin{lemma}[Zone A]\label{lem:zoneA}
Let $(q,\alpha,m)\in\mathcal D$ with $e\le\tfrac1{60}$ and $\xi\ge1$.  Then
\begin{equation}\label{eq:zoneA-target}
        R_m\ \le\
        \sqrt{2\alpha}\,B_m f\Bigl(d-\frac{e^2}{16\alpha^2(1+\rho)}\Bigr)
        \ \le\ C_m\psi(\xi,\rho).
\end{equation}
\end{lemma}

The second inequality is the dual certificate \eqref{eq:cert-II}; the content
is the first.  We may assume $R_m>0$.  All terminating decimals in this
section denote the displayed exact rational numbers.  Put
\[
        z:=(m-2)u,
        \qquad
        \eps:=\frac{e^2}{16\alpha^2(1+\rho)d}
        =\frac{e}{4(1-e)^2(1+\rho)\kappa},
        \qquad
        \sigma_e:=\left(\frac{51}{25}e\right)^6.
\]
For $e\le1/60$ and $0<\kappa\le1$ we shall use
\begin{equation}\label{eq:standing}
 \alpha\ge\frac{59}{120},\qquad
 p\le\frac{1+3e}{2}\le\frac{21}{40},\qquad
 x\ge\frac{1-e}{1+3e}\ge0.9365.
\end{equation}
\begin{equation}\label{eq:standing2}
 \frac e\alpha=\frac{2e}{1-e}\le2.04e,
 \qquad
 \ell\le\sqrt{\frac{2e}{1-e}}\le1.4262\sqrt e\le0.1842.
\end{equation}
The last two decimal comparisons follow, for example, from
$120/59<2.034$ and $(1.4262)^2>120/59$.

\begin{lemma}[Positivity threshold in Zone A]\label{lem:posA}
If $R_m>0$, then
\begin{equation}\label{eq:posA}
        \kappa(m-1)>x(1+\kappa)-\sigma_e,
        \qquad
        \kappa\ge\frac{0.93}{m-1}.
\end{equation}
\end{lemma}

\begin{proof}
From \eqref{eq:defect-form}, $R_m>0$ gives
$\tau^{m-1}<x+L^m/(p\alpha^{m-1})$.  Since $L^2\le\alpha e$ and $m\ge15$,
\[
 \frac{L^m}{p\alpha^{m-1}}
 \le x\left(\frac e\alpha\right)^{m/2}
 \le x(2.04e)^{7.5}
 =x\sigma_e(2.04e)^{1.5}.
\]
Bernoulli's inequality gives
$\tau^{m-1}\ge1-(m-1)d/\alpha$, and hence
\[
 (m-1)\frac d\alpha
 >1-x\bigl(1+\sigma_e(2.04e)^{1.5}\bigr)
 \ge\frac{d+e}{p}-\sigma_e(2.04e)^{1.5}.
\]
After multiplication by $\alpha/e$, the first term becomes $x(1+\kappa)$,
whereas
$\alpha 2.04^{1.5}\sqrt e<0.19$.  This proves the first inequality.
Since $x-\sigma_e>0.93$, it also gives the second.
\end{proof}

\begin{lemma}[Payment and defect reduction]\label{lem:reduceA}
Under the hypotheses above,
\begin{equation}\label{eq:epsA}
        \eps\le\min\left(\frac14,\,0.2781\,T\right),
        \qquad T:=me,
\end{equation}
and \eqref{eq:zoneA-target} follows from
\begin{equation}\label{eq:reducedA}
 \sqrt{1-e}\,\frac{(1-\ell)(1-y^{m-1})(1-\eps)}{1+y}
 \ge
 x^{m-2}\left[\frac{m-1}{mx}-\frac{1+1/\kappa}{m}\right]+\Lambda,
 \qquad
 \Lambda:=\frac{(\alpha e)^{m/2}}{m d\alpha p^{m-2}}.
\end{equation}
Moreover $\Lambda\le\sigma_e x^{m-2}$.
\end{lemma}

\begin{proof}
The condition $\xi\ge1$ gives $d\ge e^2/(4\alpha^2)$, hence
$\eps\le1/[4(1+\rho)]\le1/4$.  Lemma~\ref{lem:posA} also gives
\[
 \eps\le\frac{(m-1)e}{0.93\cdot4(1-e)^2}
 <0.2781\,me,
\]
because $[0.93\cdot4(59/60)^2]^{-1}<0.2781$.

Next,
\[
 B_m\ge m k_m(L)
 =\frac{m p^{m-1}(1-y^{m-1})}{p+L}.
\]
Using $\sqrt{2\alpha}=\sqrt{1-e}$, $f=\alpha(1-\ell)$ and
$p+L=p(1+y)$ gives the left side of \eqref{eq:reducedA}, multiplied by
$m d\alpha p^{m-2}$.  On the other hand, Bernoulli's inequality in
\eqref{eq:defect-form} gives
\[
 R_m\le\alpha^{m-1}\left[\frac{(m-1)d}{x}-(d+e)\right]
       +(\alpha e)^{m/2}.
\]
Division by $m d\alpha p^{m-2}$ yields \eqref{eq:reducedA}.  Finally,
\[
 \Lambda=x^{m-2}\frac{(e/\alpha)^{m/2-1}}{m\kappa}
 \le x^{m-2}\frac{(2.04e)^{6.5}}{0.93}
 \le\sigma_e x^{m-2},
\]
since $m\kappa\ge0.93$ and $\sqrt{2.04e}<0.19<0.93$.
\end{proof}

\begin{proof}[Proof of Lemma~\ref{lem:zoneA}]
Let the left and right sides of \eqref{eq:reducedA} be $\mathrm{LHS}$ and
$\mathrm{RHS}$.  Since $y\le\ell$,
$(1-\ell)/(1+\ell)\ge1-2\ell$,
$\sqrt{1-e}\ge1-0.51e$, and
$y^{m-1}\le\ell^{14}\le(2.04e)^7\le\sigma_e$, we have
\begin{equation}\label{eq:GammaA}
        \mathrm{LHS}\ge1-\Gamma,
        \qquad
        \Gamma:=2\ell+\eps+0.51e+\sigma_e.
\end{equation}
If the square bracket in \eqref{eq:reducedA} is negative, then
$\mathrm{RHS}\le\Lambda\le\sigma_e$, whereas
$\mathrm{LHS}>0.36$, and there is nothing to prove.  Assume from now on that
the bracket is nonnegative.

We have
\[
 \left[\cdots\right]
 \le\frac1x-\frac{2+1/\kappa}{m},
 \qquad
 \frac1x\le1+2.034(1+\kappa)e,
\]
and, by \eqref{eq:E-a},
\[
        x^{m-2}\le\frac1{1+z+\frac6{13}z^2}.
\]
Furthermore
\[
 u=\frac{2(1+\kappa)e}{1+e+2\kappa e}
 \ge\frac{40}{21}(1+\kappa)e>1.9047(1+\kappa)e.
\]
Consequently
\begin{align}
 2.034(1+\kappa)e&\le0.0822z,\label{eq:Aaux1}\\
 T&\le0.6058z,
 &\eps&\le0.1686z,\label{eq:Aaux2}\\
 \frac{2+1/\kappa}{m}
 &\ge\frac{11.08e}{z+7.8e}.\label{eq:Aaux3}
\end{align}
For \eqref{eq:Aaux3}, use
$(1+\kappa)(2+1/\kappa)\ge3+2\sqrt2>5.8284$ and
$1+\kappa\le2$.

Thus \eqref{eq:reducedA} follows from
\begin{equation}\label{eq:battle}
 (1-\Gamma)\left(1+z+\frac6{13}z^2\right)
 \ge1+0.0822z-\frac{11.08e}{z+7.8e}+\sigma_e,
\end{equation}
where
\[
 \Gamma\le A_0+\min\left(\frac14,0.1686z\right),
 \qquad
 A_0:=2.8527\sqrt e+0.51e+\sigma_e.
\]
The term $\sigma_e$ on the right absorbs $\Lambda$ because
$\Lambda\le\sigma_e x^{m-2}$ and
$x^{m-2}(1+z+\frac6{13}z^2)\le1$.

Let $F$ denote the left side of \eqref{eq:battle} minus the right side.  Then
\[
 F=z(0.9178-\Gamma)+\frac6{13}(1-\Gamma)z^2
   +\frac{11.08e}{z+7.8e}-\Gamma-\sigma_e.
\]
For $e\le1/60$,
\begin{equation}\label{eq:A0bounds}
        A_0\le0.3768,
        \qquad
        \sigma_e\le10^{-7}e.
\end{equation}
We check three ranges.

\smallskip\noindent
\emph{Range 1: $0<z\le1$.}
Using $\Gamma\le A_0+0.1686z$, and then $z\le1$, gives
\[
 F\ge0.3724z+\frac{11.08e}{z+7.8e}-A_0-\sigma_e;
\]
the discarded $z^2$ coefficient is at least $0.041$.  By AM--GM,
\[
 0.3724z+\frac{11.08e}{z+7.8e}
 \ge4.0626\sqrt e-2.9048e.
\]
Therefore
\[
 F\ge1.2099\sqrt e-3.4148e-2\sigma_e>0,
\]
since $3.4148\sqrt e<0.441$ and \eqref{eq:A0bounds} holds.

\smallskip\noindent
\emph{Range 2: $1\le z\le1.4827$.}
Here $0.1686z<1/4$.  Substitution of $A_0\le0.3768$ gives
\[
 F\ge q(z):=0.3724z+0.1190z^2-0.0779z^3-0.3768-\sigma_e.
\]
The derivative
$q'(z)=0.3724+0.2380z-0.2337z^2$ is positive on this interval
(it is concave and positive at both endpoints), while
$q(1)>0.036-\sigma_e>0$.

\smallskip\noindent
\emph{Range 3: $z\ge1.4827$.}
Now $\Gamma\le A_0+1/4\le0.6268$, so
\[
 F\ge0.2910z+0.1722z^2-0.6268-\sigma_e.
\]
The right side is increasing and exceeds $0.183-\sigma_e$ at $z=1.4827$.
Thus $F>0$ in every range, proving \eqref{eq:battle} and the lemma.
\end{proof}


\begin{remark}[Where Zone A is tight]\label{rmk:zoneA-tight}
In the pinch corner ($e\to0$, $T=me\to0$, $\kappa$ fixed), payment and defect
agree to leading order: both are $2\kappa m e\,p^m(1+o(1))$, and the margin is
$(4\kappa+2)e\,p^m(1+o(1))$, i.e.\ a relative margin $\sim(2+1/\kappa)/m$.
Thus the relative margin can tend to zero on scales of order $1/m$.
In particular, no proof with a uniform positive relative margin is possible;
Range~1 above must retain the $O(\sqrt e)$ and $O(1/m)$ effects rather than
discarding them.  The optional numerical scan in
\S\ref{sec:verification} is consistent with this asymptotic picture.
\end{remark}

\section{Zone B: \texorpdfstring{$e\ge1/60$, $\xi\ge1$}{e >= 1/60, xi >= 1}}\label{sec:zoneB}

\begin{lemma}[Zone B]\label{lem:zoneB}
Let $(q,\alpha,m)\in\mathcal D$ with $\tfrac1{60}\le e<\tfrac13$, $\xi\ge1$
and $R_m>0$.  Then
$R_m\le\sqrt{2\alpha}\,B_mf\bigl(d-\frac{e^2}{16\alpha^2(1+\rho)}\bigr)
\le C_m\psi(\xi,\rho)$.
\end{lemma}

As noted in \S\ref{sec:coords}, the hypotheses force $e<\tfrac{2033}{10000}$
(else $\kappa_\xi\ge\kappa_q$ and the zone is empty).  The payment and defect
reductions of Lemma~\ref{lem:reduceA} are domain-free, so, exactly as there,
it suffices to prove the reduced inequality
\begin{equation}\label{eq:reducedB}
        \Pi:=\sqrt{1-e}\,\frac{(1-\ell)(1-y^{m-1})(1-\eps)}{1+y}
        \ \ge\ x^{m-2}\Bigl[\frac{m-1}{mx}-\frac{1+1/\kappa}{m}\Bigr]+\Lambda .
\end{equation}
The new difficulty is that the worst cycle length $m$ is now an interior
maximum (roughly $m\approx1+A+1/\lambda$ below), so we first organise the
right side as a one-variable maximization.

\begin{lemma}[The $m$-maximization]\label{lem:Kmax}
Write $n=m-1\ge14$, $A=x(1+\kappa)/\kappa$, $\lambda=\ln(1/x)$, and
$K(n)=x^n(n-A)/(n+1)$, so that
$x^{m-2}\bigl[\frac{m-1}{mx}-\frac{1+1/\kappa}m\bigr]=K(n)/x^2$.  Then:
\begin{enumerate}[label=\textup{(\roman*)}]
\item If $n\le A$ the bracket is $\le0$ and \eqref{eq:reducedB} reduces to
$\Pi\ge\Lambda$.
\item $K$ is strictly log-concave on $(A,\infty)$, with stationary point
$n_*=A+s_*$, where $s_*(s_*+A+1)=(A+1)/\lambda$.  Hence
$\max_{n\ge14}K(n)=K(\max(14,n_*))$.
\item $K(14)\le x^{14}\,(14-A)_+/15$, and if $A<14$ then $n_*\le14$ if and
only if $15\,(14-A)\,\lambda\ge A+1$.
\item With $\nu:=\lambda(A+1)$ and $w:=\bigl(\sqrt{\nu^2+4\nu}-\nu\bigr)/2$,
\[
        K(n_*)=e^{-\lambda A}\,e^{-w}(1-w)\ \le\ e^{-\lambda A-2w},
\]
and the map $(a,\nu)\mapsto e^{-a-2w(\nu)}$ is decreasing in both arguments;
moreover
\[
        \lambda A\ \ge\ \nu'\ \ \text{and}\ \ \nu\ge\nu',
        \qquad
        \nu':=\Phi(e,\kappa):=\frac{2(1-e)(1+\kappa)^2e}{\kappa\,(1+e+2\kappa e)^2},
\]
so $K(n_*)\le e^{-\Phi-2w(\Phi)}$.
\end{enumerate}
\end{lemma}

\begin{proof}
(i) is clear ($\Lambda\ge$ the whole right side then).  (ii): $(\ln K)''=
-\frac1{(n-A)^2}+\frac1{(n+1)^2}<0$ since $0<n-A<n+1$; the stationarity
condition $\lambda=\frac1{n-A}-\frac1{n+1}=\frac{A+1}{(n-A)(n+1)}$ gives the
quadratic for $s_*=n_*-A$; unimodality gives the max formula.  (iii): the
first claim is $(n-A)/(n+1)$ at $n=14$; for the second,
$s\mapsto s(s+A+1)$ is increasing, so
\[
 s_*\le14-A
 \quad\Longleftrightarrow\quad
 \frac{A+1}{\lambda}=s_*(s_*+A+1)\le15(14-A),
\]
which is equivalent to $15(14-A)\lambda\ge A+1$.  (iv): from
$s_*(s_*+A+1)=(A+1)/\lambda$, multiplying by $\lambda^2$ gives
$w^2+\nu w=\nu$ with $w=\lambda s_*$, so $w=(\sqrt{\nu^2+4\nu}-\nu)/2<1$ and
$\frac{s_*}{s_*+A+1}=\frac{\lambda s_*^2}{(A+1)/\lambda\cdot\lambda}
=\frac{w^2}\nu=1-w$; hence
$K(n_*)=x^{A+s_*}\frac{s_*}{s_*+A+1}=e^{-\lambda A}e^{-w}(1-w)
\le e^{-\lambda A-2w}$ by $1-w\le e^{-w}$.  Monotonicity in $a$ is clear; in
$\nu$, $w'(\nu)=\frac12\bigl(\frac{\nu+2}{\sqrt{\nu^2+4\nu}}-1\bigr)>0$
since $(\nu+2)^2>\nu^2+4\nu$.  Finally
$\lambda A\ge uA=\frac{2(1+\kappa)e}{1+e+2\kappa e}\cdot
\frac{(1-e)(1+\kappa)}{\kappa(1+e+2\kappa e)}=\Phi$ using $\lambda\ge u=1-x$
and the chart formulas, and $\nu=\lambda(A+1)\ge\lambda A\ge\Phi$.
\end{proof}

\begin{lemma}[Zone-B battle]\label{lem:zoneB-battle}
For all $e\in[\tfrac1{60},\,\tfrac{2033}{10000}]$ and
$\kappa\in[\kappa_\xi(e),\,\min(\kappa_{\max},\kappa_q)(e)]$:
\begin{equation}\label{eq:B-battle}
        \underline\Pi(e,\kappa)\ \ge\
        \frac{\mathfrak K(e,\kappa)}{x^2}+\overline\Lambda(e),
\end{equation}
where
$\underline\Pi=\sqrt{1-e}\,\frac{(1-\bar\ell)(1-\bar\ell^{14})
(1-\bar\eps)}{1+\bar\ell}$ with $\bar\ell=\sqrt{2e/(1-e)}$,
$\bar\eps=\min\bigl(\tfrac14,\frac{e}{4(1-e)^2\kappa(1+\underline\rho)}\bigr)$,
$\underline\rho=\tfrac14\bigl(1-e^{-17.37(1+\kappa)e}\bigr)$;
$\mathfrak K=x^{14}(14-A)_+/15$ when the gate
$15(14-A)\lambda\ge A+1$, $A<14$ holds, and
$\mathfrak K=\max\bigl(x^{14}(14-A)_+/15,\;e^{-\Phi-2w(\Phi)}\bigr)$
otherwise; and
$\overline\Lambda=x^{13}\,(2e/(1-e))^{6.5}\,(1-e)^2/(15e)$.
\end{lemma}

\begin{proof}
This is a two-variable inequality between explicit elementary functions; it
is certified by exhaustive subdivision of the strip into rational boxes with
monotone endpoint bounds, in exact rational arithmetic (integer-square-root
bracketing for $\sqrt{\cdot}$, partial-sum bracketing with tail bounds for
$e^{-t}$): script \texttt{zoneB\_certifier.py}, which verifies the cover and
exits successfully.  The certificate is small: $23$ verified boxes ($3$
further boxes lie outside the admissible strip), maximum subdivision
depth $8$; each box is checkable by hand from the monotone endpoint bounds
below.
The monotonicity facts used for the box bounds are: $x$ is decreasing in $e$
and $\kappa$; $(1+\kappa)^2/\kappa$ is decreasing for $\kappa\le1$; $(1-e)e$
is increasing for $e\le\tfrac12$; $\kappa_\xi$ is increasing and
$\kappa_{\max},\kappa_q$ decreasing in $e$; $w(\cdot)$ is increasing;
$e^{-t}$ decreasing.
\end{proof}

\begin{proof}[Proof of Lemma~\ref{lem:zoneB}]
The hypotheses give $\kappa\ge\kappa_\xi(e)$ ($\xi\ge1$) and
$\kappa\le\min(\kappa_{\max},\kappa_q)$ (Lemma~\ref{lem:chart}); as computed
above the strip is empty unless $e<\tfrac{2033}{10000}$.  The quantities in
\eqref{eq:reducedB} obey
$\Pi\ge\underline\Pi$ (using $\ell\le\bar\ell$, $y\le\ell$,
$y^{m-1}\le\bar\ell^{14}$, $\eps\le\bar\eps$, the latter by
$\rho\ge(1-x^{m-1})\frac{\sqrt\alpha}{2\sqrt2f(1+x)}
\ge\tfrac14(1-e^{-14u})\ge\underline\rho$.  Here the first inequality
uses
$A_m/B_m\ge k_m(\alpha)/k_m(L)\ge(1-x^{m-1})/(1+x)$, and
$\frac{\sqrt\alpha}{2\sqrt2f(1+x)}\ge\frac1{2\sqrt2\sqrt\alpha\,(1+x)}
\ge\tfrac14$ and, because $e\le2033/10000$,
\[
 14u\ge \frac{28(1+\kappa)e}{1+3e}
 \ge17.37(1+\kappa)e
\]
(the last comparison is an exact rational inequality));
$\Lambda\le\overline\Lambda$ (the ratio of consecutive odd-$m$ values of
$\Lambda$ is $x^2(e/\alpha)\cdot\frac m{m+2}<1$, so the $m=15$ value bounds
all, and $\kappa\ge\kappa_\xi$ bounds $1/(15\kappa)$); and by
Lemma~\ref{lem:Kmax} the right side of \eqref{eq:reducedB} is at most
$\mathfrak K/x^2+\Lambda$.  Now Lemma~\ref{lem:zoneB-battle} closes
\eqref{eq:reducedB}, hence the certificate inequality, hence the lemma by
\eqref{eq:cert-II}.
\end{proof}

\section{Zone C: \texorpdfstring{$\xi\le1$}{xi <= 1}}\label{sec:zoneC}

\begin{lemma}[Zone C]\label{lem:zoneC}
Let $(q,\alpha,m)\in\mathcal D$ with $\xi\le1$ and $R_m>0$.  Then
$R_m\le C_m\psi(\xi,\rho)$.  More precisely, the proof exhibits an
admissible value of the dual parameter: in the small-$e$ range it uses
$\lambda=2\rho\xi<1$; in the moderate-$e$ range it uses either
$\lambda=2\rho_{\mathrm{lo}}\xi\le1$ or $\lambda=1$, where the explicit
lower bound $\rho_{\mathrm{lo}}\le\rho$ is defined in
Lemma~\ref{lem:threegeo}.
\end{lemma}

We prove the small-$e$ half here (Lemma~\ref{lem:zoneC-small}); the
moderate-$e$ half, which contains the Tur\'an corner, is
Lemma~\ref{lem:zoneC-mod} in the next subsection.  Note $\xi\le1$ is
equivalent to $\kappa\le\kappa_\xi:=e/(1-e)^2$, i.e.\ $d\le e^2/(1-e)^2$: the
frontier sits very close to $q$.

\subsection{Small \texorpdfstring{$e$}{e}}

\begin{lemma}[Zone C, $e\le1/60$]\label{lem:zoneC-small}
Lemma~\ref{lem:zoneC} holds when $e\le\tfrac1{60}$.
\end{lemma}

\begin{proof}
Lemma~\ref{lem:posA} applies, because its proof used only $e\le1/60$.
Thus $\kappa\ge0.93/(m-1)$.  Since
$\kappa\le\kappa_\xi=e/(1-e)^2$,
\begin{equation}\label{eq:C-mT}
        m\ge55,
        \qquad
        T=me>(m-1)e\ge0.93(1-e)^2\ge0.899.
\end{equation}

\emph{The branch $2\rho\xi>1$ is empty.}
Because $k_m(\alpha)\le k_m(L)$,
\[
 \frac{A_m}{B_m}
 \le\frac{2y^{m-2}(1+y)}{m(1-y^{m-1})}
     +\frac{1+y}{1+x}
 \le10^{-38}+0.6117.
\]
Here we used $y\le\ell\le0.1842$, $x\ge0.9365$, and $m\ge55$.
Moreover $f=\alpha(1-\ell)\ge0.8157\alpha$, so
\[
 \frac{\sqrt\alpha}{2\sqrt2f}\le0.6182.
\]
Consequently $\rho\le0.3782$ and
$2\rho\xi\le0.7564<1$.

Since $2\rho\xi<1$, the dual choice $\lambda=2\rho\xi$ gives the value in
\eqref{eq:cert-I}.  It is therefore enough to prove the slightly stronger
inequality
\begin{equation}\label{eq:Csmall-strong}
        R_m\le2\alpha^3A_m\kappa^2.
\end{equation}
Since
\[
 (m-1)u\ge1.9047(m-1)e\ge1.712,
 \qquad
 1-x^{m-1}\ge1-e^{-1.712}>0.8142,
\]
we have
\[
 A_m\ge m k_m(\alpha)
 =\frac{m p^{m-2}(1-x^{m-1})}{1+x}.
\]
Put
\[
 t:=(m-1)\log(1/\tau)-\log(1/x),
 \qquad z:=\kappa m.
\]
Then $\alpha-p\tau^{m-1}=\alpha(1-e^{-t})$ and
\begin{equation}\label{eq:Csmall-t}
 R_m\le\alpha^m\min(1,t_+)+L^m,
 \qquad
 t\le e(2.0351z-1.9047).
\end{equation}
Indeed,
$\log(1/\tau)\le d/q$, $\log(1/x)\ge1-x=u\ge1.9047e$,
$1/q\le2.0351$, and $(m-1)d\le ze$.

After division by $\alpha^3p^{m-2}$, inequality
\eqref{eq:Csmall-strong} follows if
\begin{equation}\label{eq:Csmall-normalized}
 x^{m-2}\min(1,t_+)+\ell^2y^{m-2}
 \le0.8142\alpha\kappa^2m.
\end{equation}
Now $0.8142\alpha\ge0.8142\cdot59/120>0.4003$.  Also
$\ell^2,y^2\le0.034$ and $m\ge55$, whereas
$\kappa^2m>0.864/m$; hence
\[
        \ell^2y^{m-2}<10^{-37}\kappa^2m.
\]
It therefore suffices to prove
\begin{equation}\label{eq:C-battle}
        x^{m-2}\min(1,t_+)\le0.39\kappa^2m.
\end{equation}
Since $m\ge55$,
\[
 (m-2)u\ge1.9047\frac{m-2}{m}T\ge1.8354T,
 \qquad
 x^{m-2}\le e^{-1.8354T},
 \qquad
 \kappa^2m=\frac{z^2e}{T}.
\]

If $t\le1$, then \eqref{eq:Csmall-t} and
\[
 \frac{(2.0351z-1.9047)_+}{z^2}
 \le\frac{2.0351^2}{4\cdot1.9047}<0.5437
\]
reduce \eqref{eq:C-battle} to
$Te^{-1.8354T}\le0.39/0.5437$.  This holds because
\[
        \max_{T>0}Te^{-1.8354T}
        =\frac1{1.8354\,\mathrm e}<0.2005<\frac{0.39}{0.5437}.
\]
If $t>1$, then \eqref{eq:Csmall-t} implies
$z\ge0.4753/e$.  Hence the right side of \eqref{eq:C-battle} is at least
$0.0881/(eT)$, while
\[
 eT e^{-1.8354T}\le\frac{0.2005}{60}<0.0034<0.0881.
\]
Thus \eqref{eq:C-battle}, and hence \eqref{eq:Csmall-strong}, holds.
\end{proof}


\subsection{Moderate \texorpdfstring{$e$}{e} and the Tur\'an corner}\label{subsec:zoneC-mod}

\begin{lemma}[Zone C, $e\ge1/60$]\label{lem:zoneC-mod}
Lemma~\ref{lem:zoneC} holds when $\tfrac1{60}\le e<\tfrac13$.
\end{lemma}

The proof rests on three structural facts and splits at
$\delta:=\alpha-\tfrac13=\tfrac{1-3e}6$ (note the domain constraint $q>\tfrac13$
is exactly $d<\delta$): the range $\delta\ge\tfrac1{2000}$ (i.e.\
$e\le\tfrac13-\tfrac1{1000}$) is certified by exact finite subdivision, and the
Tur\'an-corner sliver $\delta<\tfrac1{2000}$ is Lemma~\ref{lem:turan} below.

\begin{lemma}[Three-geometric defect and secant gate]\label{lem:threegeo}
On all of $\mathcal D$, with $s:=q/p$ (so $0\le y<s<x<1$):
\begin{enumerate}[label=\textup{(\roman*)}]
\item $\dfrac{R_m}{\alpha^3p^{m-2}}
=\dfrac{x^{m-2}}\alpha+\dfrac{\ell^2y^{m-2}}\alpha-\dfrac{pq}{\alpha^3}s^{m-2}$,
and $\dfrac{pq}{\alpha^3}=\dfrac1\alpha+\dfrac{\ell^2}\alpha$ identically;
consequently
\[
        R_m>0
        \iff
        x^{m-2}-s^{m-2}\ >\ \ell^2\bigl(s^{m-2}-y^{m-2}\bigr).
\]
\item (Secant gate.)  If $R_m>0$ and $m\ge15$, then
\[
        (m-2)\,\frac dq\ \ge\ (m-2)\ln\frac\alpha q\ >\ \ln(1+G_2),
        \qquad
        G_2:=\ell^2\Bigl(1-\bigl(\tfrac ys\bigr)^{13}\Bigr).
\]
\item (Payments, $m$-free-gated.)  Let
$\rho_{\mathrm{lo}}:=(1-x^{14})\frac{\sqrt\alpha}{2\sqrt2f(1+x)}$; then
$\rho\ge\rho_{\mathrm{lo}}$, and:
if $2\rho_{\mathrm{lo}}\xi\le1$,
\[
        \frac{C_m\psi(\xi,\rho)}{\alpha^3p^{m-2}}
        \ \ge\ m\cdot
        \frac{2(1-x^{14})(1-y^{m-1})\,\kappa^2\,(1+4\xi)}
        {(1+x)(1+y)(1+2\xi)}\ ;
\]
if $2\rho_{\mathrm{lo}}\xi>1$,
\[
        \frac{C_m\psi(\xi,\rho)}{\alpha^3p^{m-2}}
        \ \ge\ m\cdot
        \frac{\sqrt{2\alpha}\,(1-y^{m-1})\,f\,d\,(4\xi+1)}
        {\alpha^3\,(1+y)\,(4\xi+2)}\ .
\]
\end{enumerate}
\end{lemma}

\begin{proof}
(i) Divide $R_m=\alpha^m+L^m-pq^{m-1}$ by
$\alpha^3p^{m-2}$.  The three terms become, respectively,
\[
 \frac{x^{m-2}}{\alpha},\qquad
 \frac{\ell^2y^{m-2}}{\alpha},\qquad
 \frac{pq}{\alpha^3}s^{m-2},
\]
and the coefficient identity is $pq-\alpha^2=L^2$.

(ii) From (i),
\[
 x^{m-2}>
 s^{m-2}+\ell^2s^{m-2}\left(1-(y/s)^{m-2}\right)
 \ge s^{m-2}(1+G_2),
\]
using $m-2\ge13$ and $y/s<1$.  Hence
\[
 (m-2)\log(x/s)>\log(1+G_2),
 \qquad
 \log(x/s)=\log(\alpha/q)\le\frac{\alpha-q}{q}=\frac dq.
\]

(iii) By the mediant inequality and $k_m(\alpha)\le k_m(L)$,
\[
        \frac{A_m}{B_m}\ \ge\ \frac{k_m(\alpha)}{k_m(L)}
        =\frac{(p^{m-1}-\alpha^{m-1})(p+L)}{(p^{m-1}-L^{m-1})(p+\alpha)}
        \ \ge\ (1-x^{m-1})\,\frac{p}{p+\alpha}
        \ \ge\ \frac{1-x^{14}}{1+x},
\]
giving $\rho\ge\rho_{\mathrm{lo}}$.
If $2\rho_{\mathrm{lo}}\xi\le1$, take $\lambda=2\rho_{\mathrm{lo}}\xi\in[0,1]$
in the dual form of $\psi$: since the dual value
$\lambda\xi-\frac{\lambda^2}{4(\rho+\lambda)}$ is increasing in $\rho$,
\[
        C_m\psi\ \ge\
        C_m\Bigl(2\rho_{\mathrm{lo}}\xi^2
        -\frac{4\rho_{\mathrm{lo}}^2\xi^2}{4(\rho_{\mathrm{lo}}
        +2\rho_{\mathrm{lo}}\xi)}\Bigr)
        =C_m\rho_{\mathrm{lo}}\xi^2\,\frac{1+4\xi}{1+2\xi}
        =\frac{2\alpha^3B_m(1-x^{14})\kappa^2}{1+x}\cdot\frac{1+4\xi}{1+2\xi},
\]
by the exact algebra $C_m\xi^2=4\alpha^2B_mf\sqrt{2\alpha}\,d^2/e^2$ and the
definition of $\rho_{\mathrm{lo}}$; then insert
$B_m\ge mp^{m-2}(1-y^{m-1})/(1+y)$.  If $2\rho_{\mathrm{lo}}\xi>1$ then
$2\rho\xi>1$, so $\xi>\xi_c(\rho)$ and
$\psi=\xi-\frac1{4(1+\rho)}\ge\xi\frac{4\xi+1}{4\xi+2}$
(the last step is equivalent to $4\xi\rho\ge2$); multiply by
$C_m$, use $C_m\xi=\sqrt{2\alpha}B_mfd$ and the same $B_m$ bound.
\end{proof}

\begin{proof}[Proof of Lemma~\ref{lem:zoneC-mod} for $e\le\tfrac13-\tfrac1{1000}$]
By Lemma~\ref{lem:threegeo}, it suffices that for every box of a finite cover
of $\{(e,\kappa)\}$ and every $m\ge15$, the (upper-bounded) three-geometric
defect not exceed $m$ times the (lower-bounded) applicable payment
coefficient; positive defect only needs checking for
$m\ge M_+:=3+\lfloor q(G_2-G_2^2/2)/d\rfloor$ (the strict secant
gate with $\ln(1+t)\ge t-t^2/2$), and all sufficiently large $m$ pass at once because
both positive defect terms decrease geometrically while the payment grows
linearly.  The script \texttt{zoneC\_certifier.py} certifies exactly this, in
exact rational arithmetic, over $e\in[\tfrac1{60},\tfrac13-\tfrac1{1000}]$,
$\kappa\in(0,1]$: $2997$ boxes verified (of which $5$ by the analytic
$\kappa\to0$ bottom-out, which combines the gate $m-2>C_0/\kappa$ with the
monotonicity of $e^{-a/\kappa}/\kappa$), maximum subdivision depth $28$,
largest explicitly checked cycle length $m=1716$.  The interval
construction uses the lower endpoint
\[
        f_{\mathrm{lo}}
        =\max\{\alpha_{\mathrm{lo}}-L_{\mathrm{up}},
                d_{\mathrm{lo}}+\delta_{\mathrm{lo}}\},
\]
whose validity is exactly Lemma~\ref{lem:frontier-gap}; this is the bound that
was missing from the supplied manuscript.
\end{proof}

\subsection{The Tur\'an-corner sliver}

\begin{lemma}[Tur\'an corner]\label{lem:turan}
Lemma~\ref{lem:zoneC} holds when
$0<\delta=\tfrac{1-3e}6<\tfrac1{2000}$
(equivalently $e>\tfrac13-\tfrac1{1000}$).
\end{lemma}

\begin{proof}
Write
\[
 \alpha=\frac13+\delta,\qquad e=\frac13-2\delta,\qquad
 q=\frac13+\delta-d,\qquad p=\frac23-\delta+d ,
\]
where $0<d<\delta\le1/2000$.  Since
\[
 \frac{\alpha}{e}\le\frac{2003}{1994},
 \qquad
 \xi=4\left(\frac{\alpha}{e}\right)^2d<4.04\,d<1,
\]
we are indeed in Zone~C.

We first collect bounds that will be used repeatedly.  Put
$S:=q^2-L^2=(q-L)(q+L)$.  Direct expansion gives
\begin{equation}\label{eq:corner-S}
 S=\delta+3\delta^2-\left(\frac13+4\delta\right)d+2d^2 .
\end{equation}
Writing $d=t\delta$, $0\le t\le1$, shows
\[
 \frac{S}{\delta}
 =1-\frac t3+\delta(3-4t+2t^2),
 \qquad
 \frac23\le\frac{S}{\delta}\le1+3\delta\le\frac{2003}{2000}.
\]
Moreover
\[
 L^2-e^2
 =\delta-\frac d3-6\delta^2+2\delta d-d^2
 \ge \frac23\delta-5\delta^2>0,
\]
because the displayed expression is decreasing in $d\in[0,\delta]$.
Thus $e<L<q$, and consequently
\begin{equation}\label{eq:corner-Q}
 \frac{1997}{3000}
 \le q+L\le\frac{2003}{3000}.
\end{equation}
Also
\[
 \frac12\le x=\frac{\alpha}{p}
 \le\frac12+2.26\,\delta\le0.50113,
 \qquad y<\frac12,
 \qquad x^{14}<6.4\cdot10^{-5}.
\]
Here $2\alpha-p=3\delta-d>0$, while
$x-\tfrac12=(3\delta-d)/(2p)\le2.26\delta$.

The identities
\[
 q-L=\frac{S}{q+L},
 \qquad
 f=\alpha-L=\frac{\alpha^2-L^2}{\alpha+L}
   =\frac{3\alpha\delta+d(d+e)}{\alpha+L}
\]
and \eqref{eq:corner-S}--\eqref{eq:corner-Q} imply
\begin{equation}\label{eq:corner-standing}
 \frac{2000}{2003}\,\delta\le q-L\le1.505\,\delta,
 \qquad
 \frac{3000}{2003}\,\delta\le f\le2.005\,\delta,
 \qquad
 1-\ell=\frac f\alpha\le6.02\,\delta .
\end{equation}
For the upper bound on $f$ we used $L\ge e$, and hence
$\alpha+L\ge2/3-\delta$.

\emph{Step 1: positive defect forces a visible frontier gap.}
Set
\[
 t_0:=1-\frac ys=\frac{q-L}{q}.
\]
The preceding estimates give
\[
 2.99\,\delta\le t_0\le4.52\,\delta,
 \qquad
 \ell^2\ge1-12.1\,\delta.
\]
Taylor's formula with nonnegative third-order remainder gives
$1-(1-t)^{13}\ge13t-78t^2$ for $0\le t\le1$.  Therefore
\[
 G_2=\ell^2\bigl(1-(1-t_0)^{13}\bigr)
 \ge(1-12.1\delta)
       \bigl(13\cdot2.99\delta-78(4.52\delta)^2\bigr)
 >37.8\,\delta.
\]
(The quotient of the middle expression by $\delta$ decreases on
$[0,1/2000]$, and its endpoint value is greater than $37.84$.)
On the other hand $G_2\le13t_0\le58.8\delta\le0.0294$.  Hence
\[
 \log(1+G_2)\ge G_2\left(1-\frac{G_2}{2}\right)>37.2\,\delta.
\]
Lemma~\ref{lem:threegeo}(ii), together with $q>1/3$, now yields
\begin{equation}\label{eq:corner-forcing}
 R_m>0\quad\Longrightarrow\quad
 (m-2)d>12.4\,\delta .
\end{equation}

\emph{Step 2: an upper bound on the normalized defect.}
Put $n=m-2$.  From Lemma~\ref{lem:threegeo}(i), the mean-value theorem on
$[s,x]$ and $[y,s]$, and $y=\ell x$, we obtain
\[
 \frac{\alpha R_m}{\alpha^3p^{m-2}}
 \le n(x-s)x^{n-1}-\ell^2n(s-y)y^{n-1}
 =\frac{n x^{m-3}}p\,N,
\]
where
\[
 N:=d-\frac{S}{q+L}\,\ell^{m-1}.
\]
Using \eqref{eq:corner-S}, dropping the nonpositive
$-2d^2\ell^{m-1}/(q+L)$ term, using
$\ell^{m-1}\le1$ in the positive $d$ term, and using
\[
 \ell^{m-1}\ge1-(m-1)(1-\ell)
             \ge1-6.02(m-1)\delta
\]
in the negative $\delta$ term, gives
\[
 N\le
 d\left(1+\frac{2\alpha-e}{q+L}\right)
 -\frac{3\alpha\delta}{q+L}
 +\frac{3\alpha\cdot6.02(m-1)\delta^2}{q+L}.
\]
Since $d\le\delta$,
\[
 d\left(1+\frac{2\alpha-e}{q+L}\right)
 -\frac{3\alpha\delta}{q+L}
 \le\delta\,\frac{L-e-d}{q+L}
 \le\frac{3\delta^2}{q+L}<4.51\delta^2.
\]
The last term is less than $9.06m\delta^2$, by
$\alpha\le2003/6000$ and \eqref{eq:corner-Q}.  Thus
$N<9.5m\delta^2$.  Finally
$(p\alpha)^{-1}<4.51$, and therefore
\begin{equation}\label{eq:corner-defect}
 \frac{R_m}{\alpha^3p^{m-2}}
 <43\,m(m-2)\delta^2x^{m-3}.
\end{equation}

\emph{Step 3: the two payments.}
Lemma~\ref{lem:threegeo}(iii), together with
$\kappa=d/e\ge3d$, $y<1/2$, \eqref{eq:corner-standing}, and the bounds above,
gives
\begin{align}
 2\rho_{\mathrm{lo}}\xi\le1
 &\quad\Longrightarrow\quad
 \frac{C_m\psi(\xi,\rho)}{\alpha^3p^{m-2}}
 >7.9\,d^2m,                                      \label{eq:corner-payI}\\
 2\rho_{\mathrm{lo}}\xi>1
 &\quad\Longrightarrow\quad
 \frac{C_m\psi(\xi,\rho)}{\alpha^3p^{m-2}}
 >10.9\,\delta d\,m.                              \label{eq:corner-payII}
\end{align}
For completeness, the first numerical coefficient follows from
\[
 \frac{2(1-6.4\cdot10^{-5})(1-2^{-14})\,9}
      {1.50113\cdot1.5}>7.9.
\]
For the second, use
$\sqrt{2\alpha}>0.816$, $f>1.497\delta$,
$\alpha^3<0.0373$, $(4\xi+1)/(4\xi+2)\ge1/2$, and
\[
 \frac{0.816(1-2^{-14})\,1.497}{2\cdot0.0373\cdot1.5}>10.9.
\]

\emph{Step 4: comparison.}
Assume $R_m>0$.  In the first branch, \eqref{eq:corner-forcing} and
\eqref{eq:corner-payI} reduce the desired comparison with
\eqref{eq:corner-defect} to
\[
 43(m-2)^3x^{m-3}<7.9\cdot12.4^2.
\]
In the second branch, \eqref{eq:corner-payII} reduces it to
\[
 43(m-2)^2x^{m-3}<10.9\cdot12.4.
\]
For $j=2,3$, the sequence $(m-2)^j x^{m-3}$ decreases for $m\ge15$, since
\[
 x\left(\frac{m-1}{m-2}\right)^3
 \le0.50113\left(\frac{14}{13}\right)^3<1.
\]
It therefore suffices to check $m=15$, where
\[
 43\cdot13^3\cdot0.50113^{12}<24
 <7.9\cdot12.4^2,
 \qquad
 43\cdot13^2\cdot0.50113^{12}<1.83
 <10.9\cdot12.4.
\]
Thus the applicable payment dominates the defect in either branch.
\end{proof}

\begin{proof}[Proof of Lemma~\ref{lem:zoneC-mod}]
Combine the certified range $e\le\tfrac13-\tfrac1{1000}$ with
Lemma~\ref{lem:turan}.
\end{proof}

\section{Assembly}\label{sec:assembly}

\begin{proof}[Proof of Theorem~\ref{thm:scalar}]
Let $(q,\alpha,m)\in\mathcal D$.  If $R_m\le0$ there is nothing to prove, since
$\psi\ge0$ and $C_m>0$ ($B_m>0$ always, and $f>0$ because
$L^2\le\alpha e<\alpha^2$: indeed $e<\tfrac13<\tfrac{1-e}2=\alpha$).
So assume $R_m>0$.  If $\xi\le1$, apply
Lemma~\ref{lem:zoneC}.  If $\xi\ge1$ and $e\le\tfrac1{60}$, apply
Lemma~\ref{lem:zoneA}; if $\xi\ge1$ and $e\ge\tfrac1{60}$, apply
Lemma~\ref{lem:zoneB}.  In each case the certified quantity is a value of the
dual functional $\lambda\mapsto C_m(\lambda\xi-\frac{\lambda^2}{4(\rho+\lambda)})$
at an admissible $\lambda\in[0,1]$, hence $\le C_m\psi(\xi,\rho)$ by
Proposition~\ref{prop:psi-forms}.
\end{proof}

\begin{proof}[Proof of Corollary~\ref{cor:region2}]
Let $W$ have density $p\in(\tfrac12,\tfrac23)$ and let $m\ge3$ be odd.  For
$m\le13$ the conjecture is proved in \texttt{earlier\_consolidated\_note.tex} at all densities.
Let $m\ge15$.  If $\lambda_{\max}(A)\le q$, the bound follows from
Lemma~\ref{lem:no-frontier}.
Otherwise, by Lemmas~\ref{lem:one-frontier} and~\ref{lem:forced-variance}, $A$
has a unique frontier eigenvalue $\alpha\in(q,r(q)]$, so $(q,\alpha,m)\in\mathcal D$
and Theorem~\ref{thm:scalar} gives $R_m\le C_m\psi(\xi,\rho)$.  By
Theorem~\ref{thm:huber} and Proposition~\ref{prop:master-defect},
\[
        t(C_m,W)-\bigl(p^m-pq^{m-1}\bigr)
        \ \ge\ -R_m+A_mc^2+B_m\|g_s\|_2^2
        \ \ge\ -R_m+C_m\psi(\xi,\rho)\ \ge\ 0. \qedhere
\]
\end{proof}

\begin{proof}[Proof of Corollary~\ref{cor:full}]
For $p\le\tfrac12$ the right-hand side of \eqref{eq:conj} is nonpositive.  For
$p\ge\tfrac23$ this is the high-density theorem of \texttt{earlier\_consolidated\_note.tex}.
For $\tfrac12<p<\tfrac23$ this is Corollary~\ref{cor:region2}.
\end{proof}

\section{Exact certificate protocol}\label{sec:certificate-protocol}

This section records precisely what the two finite verifiers do.  It is part
of the proof specification: a replacement implementation is valid if it
checks the same interval inequalities with outward rounding.

The programs certify the relative interiors of their strips.  A box that lies
entirely on an overlap or an excluded boundary may be discarded by a
one-sided domain test: the overlap $\xi=1$ is already covered by Zone C, the
boundary $q=1/3$ is outside the stated domain, and the admissible ceiling
$\alpha=r(q)$ follows from the certified interior by continuity of all scalar
quantities.  Thus no admissible endpoint is lost by the strict skip tests.

\subsection{Arithmetic primitives}

Both certifiers use Python's \texttt{fractions.Fraction}; hence every field
operation and every comparison is exact.  Square roots are enclosed by
integer square roots: for a fixed decimal scale $S$,
\[
 \frac{\lfloor\sqrt{\lfloor tS^2\rfloor}\rfloor}{S}
 \le\sqrt t\le
 \frac{\lfloor\sqrt{\lfloor tS^2\rfloor}\rfloor+1}{S}.
\]
For $t\ge0$, an upper bound on $e^{-t}$ is obtained by first rounding $t$
down rationally and then using
\[
        e^{-t}\le
        \left(\sum_{j=0}^{N}\frac{t^j}{j!}\right)^{-1}.
\]
Where a lower bound is requested by the diagnostic checker, the omitted
positive tail of the exponential series is bounded by a geometric tail.
Integer powers in the Zone-C verifier are evaluated by binary powering with
a directed rational rounding after every multiplication.  The only
floating-point conversions in the exact certifiers occur in error messages
or printed summaries; no branch that returns ``verified'' depends on them.

\subsection{Zone B}

The verifier starts from
\[
 \frac1{60}\le e\le\frac{2033}{10000},
 \qquad
 \kappa_\xi(e)\le\kappa\le
 \min\{\kappa_{\max}(e),\kappa_q(e)\}.
\]
A rational rectangle is discarded only if monotonic endpoint tests show
that it lies wholly outside this strip.  On every remaining rectangle the
program forms outward bounds for $x,A,\lambda,\bar\ell,\rho,\varepsilon,
\Phi$, the endpoint value $K(14)$, and the interior maximum bound
$e^{-\Phi-2w(\Phi)}$.  It then checks exactly the box version of
\eqref{eq:B-battle}.  A failed box is bisected in the coordinate with the
larger relative width.  The run terminates with $23$ verified leaves,
$3$ discarded leaves and maximum depth $8$; reaching the depth cap with an
unverified box causes a nonzero exit.

\subsection{Zone C on the compact middle range}

Here the initial domain is
\[
 \frac1{60}\le e\le\frac13-\frac1{1000},
 \qquad
 0<\kappa\le\min\{\kappa_\xi(e),\kappa_{\max}(e),\kappa_q(e)\}.
\]
For each rational box the verifier encloses
$\alpha,d,q,p,x,s,L,y,\ell,\xi$ and the two payment coefficients in
Lemma~\ref{lem:threegeo}(iii).  In particular it uses
\[
 f_{\rm lo}=
 \max\{\alpha_{\rm lo}-L_{\rm up},\,d_{\rm lo}+\delta_{\rm lo}\},
\]
which is valid by Lemma~\ref{lem:frontier-gap}.  The secant gate says that
positive defect is possible only for
\[
 m\ge3+\left\lfloor
 \frac{q_{\rm lo}\bigl(G_{2,\rm lo}-G_{2,\rm lo}^2/2\bigr)}
      {d_{\rm up}}\right\rfloor .
\]
Starting there, the program checks each integer $m$ with an outward upper
bound on the three-geometric defect.  Once the sum of the two positive
geometric terms is below the linearly growing payment, all later $m$ are
certified at once.  Boxes touching $\kappa=0$ are closed by the analytic
bottom-out estimate
\[
 m-2>\frac{C_0}{\kappa},\qquad
 \frac{R_m}{\alpha^3p^{m-2}}
 \le\frac1{\alpha_{\rm lo}}e^{-a/\kappa},
\]
together with monotonicity of $e^{-a/\kappa}/\kappa$ on the verified range.
The completed run has $2997$ verified leaves, including $5$ bottom-out
leaves, $87$ discarded leaves, maximum depth $28$, and largest explicitly
tested exponent $m=1716$.  As in Zone B, any unresolved leaf causes a
nonzero exit.

\section{Verification}\label{sec:verification}

Following the standing practice of this project, every claim above carries
one of three labels, and the accompanying scripts let a reader re-check each
label mechanically.

\begin{enumerate}[label=\textup{(\arabic*)}]
\item \emph{Analytic.}  The one-frontier reduction, including the
no-frontier case (\S\ref{sec:reduction}), the coordinate and defect identities
(\S\ref{sec:coords}), Zone A, the $m$-maximization
Lemma~\ref{lem:Kmax}, the Zone-C small-$e$ lemma, the Tur\'an-corner lemma, and
the assembly are complete pencil-and-paper proofs.  Every terminating decimal
used in those arguments is checked as an exact rational inequality by the
first layer of the companion checker.
\item \emph{Exact certificate.}  The Zone-B battle
(Lemma~\ref{lem:zoneB-battle}) and the moderate-$e$ part of Zone C
(Lemma~\ref{lem:zoneC-mod}, range $e\le\tfrac13-\tfrac1{1000}$) are proved by
finite exhaustive box subdivisions in \emph{exact rational arithmetic}
(\texttt{zoneB\_certifier.py}: $23$ boxes; \texttt{zoneC\_certifier.py}:
$2997$ boxes with a per-box finite $m$-scan plus gate and tail lemmas); no floating point enters the certified
statements.  The transcendentals $\sqrt t$ and $e^{-t}$ are bracketed by
integer square roots and by exponential-series partial sums with explicit
tail bounds.
\item \emph{Numerical (safety net only; no theorem depends on it).}
\texttt{regionII\_scalar\_checker.py} first checks the algebraic identities
and all directed constants exactly (using \texttt{Fraction} and SymPy), then
tests every analytic chain on adversarial high-precision grids.  With
\texttt{--scan}, the recorded reproducible run examined $1359$
positive-defect parameter triples, through $m=1{,}000{,}001$ on the prescribed grid.
It found no target violation and no point missed by both dual certificates;
the smallest relative payment margin in that run was
$4.638\cdot10^{-5}$, at
$(e,\kappa,m)=(10^{-10},\,0.59999999988,\,100001)$.  These data are useful for
transcription control only and are not invoked by any proof.
\end{enumerate}

\subsection*{How to audit}

\begin{enumerate}[label=\textup{(\alph*)}]
\item Audit \S\ref{sec:reduction} independently from the block
decomposition onward.  The points most easily missed are
Remarks~\ref{rmk:budget-sign}, \ref{rmk:knegative}, the no-frontier
Lemma~\ref{lem:no-frontier}, and the $K=0$ case in
Theorem~\ref{thm:huber}.
\item Check the chart Lemma~\ref{lem:chart} by direct computation (or run
the exact identity layer of the checker).
\item Zone A: the only nontrivial verifications are the elementary
inequality \eqref{eq:E-a}, the constants in \eqref{eq:standing}, and the
three ranges of the battle \eqref{eq:battle}; each range is a one-variable
calculus argument spelled out in full.
\item Zone B: verify Lemma~\ref{lem:Kmax} (five lines of calculus), then
run \texttt{zoneB\_certifier.py}.
\item Zone C: the small-$e$ case reuses the Zone-A machinery plus the
battle \eqref{eq:C-battle}; the moderate-$e$ case is
Lemma~\ref{lem:threegeo} (a page of identities) plus
\texttt{zoneC\_certifier.py}; the Tur\'an-corner Lemma~\ref{lem:turan} is
self-contained, and every displayed constant of its proof is re-verified
as an exact rational comparison by the checker.
\item Run \texttt{python3 regionII\_scalar\_checker.py --scan} for the
numerical safety net.
\end{enumerate}

\end{document}
